\documentclass[12pt, letterpaper]{article}

% ==========================================
%              DOCUMENT PACKAGES
% ==========================================

% --- Core Language & Encoding ---
\usepackage[english]{babel} % Sets the language and regional settings for English documents
\usepackage[utf8]{inputenc} % Allows the use of UTF-8 characters in the document
\usepackage[T1]{fontenc} % Specifies the font encoding
\usepackage[margin=1.5cm]{geometry} % Sets the page margins

% --- Colors & Graphics ---
\usepackage[table,dvipsnames]{xcolor} % Advanced color management
% \usepackage{color} % Basic package for text and background colors
\usepackage{graphicx} % Enhanced support for including external images and scaling them
\usepackage{adjustbox} % Adjusts the size, scaling, and alignment of boxes, images, and tables
% \usepackage{background} % Allows adding background images or watermarks to pages
\usepackage{float} % Improves the interface for floating objects and adds the [H] placement
\usepackage{pdflscape} % Makes landscape pages display correctly on screen in PDF viewers
\usepackage{pdfpages} % Allows the inclusion of external multi-page PDF documents

% --- Math, Science & Engineering ---
\usepackage{amsfonts} % Provides extended mathematical fonts
\usepackage{amsmath} % Provides advanced mathematical typesetting and environments
\usepackage{amssymb} % Provides extended mathematical symbols
% \usepackage{amsmath, amssymb} % (Redundant combination line omitted for stability)
\usepackage{amsthm} % Facilitates the creation of theorem-like environments
\usepackage{mathtools} % Enhances amsmath with additional math tools and bug fixes
\usepackage{thmtools} % Extensions and customization options for theorem environments
\usepackage{nicematrix} % Enhances matrices and tabulars using PGF/TikZ
\usepackage{xy} % Typesets complex graphs and diagrams
\usepackage{circuitikz} % Allows drawing of electrical networks and circuits

% --- Typography & Fonts ---
\usepackage{avant} % Sets the default sans-serif font to Avant Garde
\usepackage{cantarell} % Sets the default sans-serif font to Cantarell
\usepackage{charter} % Sets the default roman font to Charter
\usepackage{cmbright} % Sets the font family to Computer Modern Bright 
\usepackage{helvet} % Sets the default sans-serif font to Helvetica
\usepackage{inconsolata} % Sets the default typewriter (monospace) font to Inconsolata
\usepackage{iwona} % Sets the font family to Iwona 
\usepackage{kurier} % Sets the font family to Kurier 
\usepackage{lmodern} % Loads the Latin Modern fonts 
\usepackage{textcomp} % Provides additional text symbols
\usepackage{ulem} % Provides various types of text underlining and strikethroughs

% --- Tables & Lists ---
\usepackage{array} % Extends the array and tabular environments
\usepackage{arydshln} % Allows dashed lines in arrays and tables
\usepackage{bigstrut} % Adds custom vertical space in table rows
\usepackage{booktabs} % Enhances the quality of tables with better horizontal rules
% \usepackage{colortbl} % Allows coloring of rows, columns, and individual cells in tables
\usepackage{diagbox} % Creates table cells with diagonal lines 
\usepackage{enumitem} % Provides advanced control over list environments
\usepackage{longtable} % Allows tables to break across multiple pages
\usepackage{multirow} % Allows individual table cells to span across multiple rows
\usepackage{tabularx} % Creates tables with dynamically adjustable-width columns
\usepackage{pgfplotstable} % Reads and cleanly formats numerical tables directly from CSV

% --- Formatting, Layout & UI ---
\usepackage{blindtext} % Generates dummy text for testing document layouts
\usepackage{caption} % Customizes captions in floating environments
\usepackage{comment} % Allows you to block out large sections of text
\usepackage{emptypage} % Automatically removes headers and footers on empty pages
\usepackage{fancyhdr} % Allows complete customization of headers and footers
\usepackage{fancyvrb} % Enhanced verbatim environments for writing raw code
\usepackage{framed} % Creates framed, shaded, or differently highlighted text regions
\usepackage{fullpage} % Quickly sets all page margins to 1 inch
\usepackage{lipsum} % Generates Lorem Ipsum dummy text
\usepackage{mdframed} % Creates highly customizable framed environments
\usepackage{multicol} % Allows intermixing single and multiple text columns
\usepackage{parskip} % Replaces standard paragraph indentation with vertical spacing
\usepackage{setspace} % Controls line spacing (single, one-and-a-half, double)
\usepackage{subcaption} % Supports sub-figures and sub-tables with independent captions
% \usepackage{subfig} % Older package for sub-figures (subcaption is generally preferred now)
\usepackage{tcolorbox} % Creates heavily customized colored and framed text boxes
\usepackage{titlesec} % Customizes the formatting of section, subsection, and chapter headings
\usepackage{titling} % Provides advanced control over the typesetting of the document title

% --- Bibliographies & References ---
\usepackage{biblatex} % Advanced, modern bibliographic processing and formatting
\usepackage{csquotes} % Context-sensitive quotation facilities
\usepackage{makeidx} % Supports the generation of a document index
% \usepackage{natbib} % Flexible bibliography support, particularly for author-year citation styles

% --- TikZ & Plotting ---
\usepackage{tikz} % Advanced package for creating vector graphics and diagrams programmatically
\usepackage{pgf-pie} % Draws pie charts programmatically using PGF/TikZ
\usepackage{pgfplots} % Creates high-quality mathematical plots, charts, and graphs
% \usepackage{fancytikzposter} % Used to create large format academic posters using TikZ

% --- URLs & Hyperlinks ---
\usepackage{url} % Formats and properly hyphenates long URLs
\usepackage{breakurl} % Allows line breaks in URLs, especially when compiled with latex -> dvips
\usepackage{hyperref} % Creates cross-referencing, clickable links, and PDF bookmarks

% --- Code Formatting ---
\usepackage{listings} % Typesets programming source code with formatting and syntax highlighting


% ==========================================
%           DOCUMENT CONFIGURATION
% ==========================================

% Custom color definitions
\definecolor{mygray}{gray}{0.9}

% Hyperlink configuration
\hypersetup{
    colorlinks=true,
    linkcolor=blue,
    urlcolor=blue
}

% Configuración adicional para breaks en URLs
\def\UrlBreaks{\do\/\do-\do_\do.\do=\do?\do\&}
\expandafter\def\expandafter\UrlBreaks\expandafter{\UrlBreaks\do\a\do\b\do\c\do\d\do\e\do\f\do\g\do\h\do\i\do\j\do\k\do\l\do\m\do\n\do\o\do\p\do\q\do\r\do\s\do\t\do\u\do\v\do\w\do\x\do\y\do\z}

%Uncomment the following lines to use a sans-serif font
\renewcommand{\familydefault}{\sfdefault}

% TikZ libraries
\usetikzlibrary{shapes.geometric, positioning, decorations.pathreplacing}

\graphicspath{ {./img/} } % Sets the path for the graphics

\newlist{biblio}{enumerate}{1}
\setlist[biblio,1]{label=[\arabic*]}

\makeindex

% \tiny
% \scriptsize
% \footnotesize
% \small
% \normalsize (default size)
% \large
% \Large
% \LARGE
% \huge
% \Huge


% ==========================================
%         GLOBAL LISTINGS SETTINGS
% ==========================================
\lstset{
	basicstyle=\ttfamily\small,
	numbers=left,
	numberstyle=\tiny\color{gray},
	stepnumber=1,
	numbersep=10pt,
	backgroundcolor=\color{white},
	showspaces=false,
	showstringspaces=false,
	breaklines=true,
	frame=single,
	tabsize=2, 
    captionpos=b, 
    lineskip=-1pt,
	extendedchars=true,
	literate={á}{{\'a}}1 {é}{{\'e}}1 {í}{{\'i}}1 {ó}{{\'o}}1 {ú}{{\'u}}1 {Á}{{\'A}}1 {É}{{\'E}}1 {Í}{{\'I}}1 {Ó}{{\'O}}1 {Ú}{{\'U}}1 {ñ}{{\~n}}1 {Ñ}{{\~N}}1 {¿}{{\textquestiondown}}1 {¡}{{\textexclamdown}}1
}

% ==========================================
%            LANGUAGE: SQL
% ==========================================
\lstdefinestyle{sqlstyle}{
	language=SQL,
	keywordstyle=\color{blue}\bfseries,
	commentstyle=\color[rgb]{0,0.5,0}\itshape,
	stringstyle=\color{red}
}

% ==========================================
%            LANGUAGE: VHDL
% ==========================================
\lstdefinelanguage{VHDL}{
	morekeywords={entity, architecture, signal, port, in, out, inout, begin, end, process, if, then, else, elsif, case, when, others, for, while, loop, wait, library, use, std_logic, std_logic_vector, integer, boolean, rising_edge, falling_edge, component, generic, map, downto, to, and, or, not, xor, nand, nor, xnor},
	morecomment=[l]{--},
	morestring=[b]",
	sensitive=false,
	alsoletter={_}
}

\lstdefinestyle{vhdl}{
	language=VHDL,
	keywordstyle=\color{blue}\bfseries,
	commentstyle=\color[rgb]{0,0.6,0}\itshape,
	stringstyle=\color{red}
}

% ==========================================
%              LANGUAGE: C
% ==========================================
\lstdefinestyle{cstyle}{
	language=C,
	keywordstyle=\color{blue}\bfseries,
	commentstyle=\color[rgb]{0,0.5,0}\itshape,
	stringstyle=\color{red},
	directivestyle=\color{purple},
    tabsize=4 % Overrides the global tabsize of 2 for C code specifically
}

% ==========================================
%             LANGUAGE: C++
% ==========================================
\lstdefinestyle{cppstyle}{
	language=C++,
	keywordstyle=\color{blue}\bfseries,
	commentstyle=\color[rgb]{0,0.5,0}\itshape,
	stringstyle=\color{red},
	directivestyle=\color{purple},
    tabsize=4
}

% ==========================================
%           LANGUAGE: PYTHON
% ==========================================
\lstdefinestyle{pythonstyle}{
	language=Python,
	keywordstyle=\color{blue}\bfseries,
	commentstyle=\color[rgb]{0,0.5,0}\itshape,
	stringstyle=\color{red},
    % Python relies heavily on indentation, so keeping tabsize aligned is important
    tabsize=4
}

% ==========================================
%             LANGUAGE: BASH
% ==========================================
\lstdefinestyle{bashstyle}{
	language=bash,
	keywordstyle=\color{blue}\bfseries,
	commentstyle=\color[rgb]{0,0.5,0}\itshape,
	stringstyle=\color{red},
    % Bash variables often use $, so we can highlight them dynamically if needed
    moredelim=[s][\color{orange}]{\$\{}{\}}
}

% ==========================================
%           LANGUAGE: ASSEMBLY
% ==========================================
\lstdefinestyle{asmstyle}{
	language=[x86masm]Assembler,
	keywordstyle=\color{blue}\bfseries,
	commentstyle=\color[rgb]{0,0.5,0}\itshape,
	stringstyle=\color{red},
    % Adding common ARM/Microcontroller instructions for better highlighting
    morekeywords={ldr, str, mov, cmp, add, sub, b, bl, bx, push, pop, svc, eor, and, orr, lsl, lsr}
}

% % Configure header
\setlength{\headheight}{28pt} % Adjust header height for two-line header
\pagestyle{fancy}
\fancyhf{} % Clear all header and footer fields
\lhead{\textbf{Nombre:} Ugartechea González Luis Antonio -- 320223822 \\ \textbf{Asignatura:} <asignatura>, \textbf{Grupo:} <grupo>}
\renewcommand{\headrulewidth}{0.4pt}

\begin{document}

\newgeometry{margin=1.5cm}
\begin{titlepage}
    \noindent
    \renewcommand{\arraystretch}{1.5}
    % Definimos la tabla principal con dos columnas exactamente iguales (50/50)
    \begin{tabularx}{\textwidth}{|>{\centering\arraybackslash}X|>{\centering\arraybackslash}X|}
        \hline
        % Fila superior: Combinamos las dos columnas e insertamos una sub-tabla
        \multicolumn{2}{|@{}c@{}|}{%
            \begin{tabularx}{\dimexpr\textwidth-2\arrayrulewidth\relax}{m{3.5cm}|>{\centering\arraybackslash}X}
                \centering \vspace{0.2cm} \includegraphics[width=2cm,height=2.5cm]{escudoFI_color.png} & \large \textbf{Fundamentos de Sistemas Embebidos}
            \end{tabularx}%
        } \\
        \hline
        % Fila inferior: Alineada a la mitad exacta
        \large Facultad de Ingeniería & \large UNAM \\
        \hline
    \end{tabularx}

    \vspace{0.5cm}

    \begin{center}
        {\fontsize{35}{48}\selectfont Laboratorio de Fundamentos de} \\
        {\fontsize{35}{48}\selectfont Sistemas Embebidos\par}
        % Línea divisoria azul (se puede ajustar el tono de azul)
        \textcolor{MidnightBlue!70}{\rule{\textwidth}{1pt}}
    \end{center}

    \vspace{2cm}

    \begin{center}
        % Aumenta el espaciado entre filas para que coincida con el formato
        % \renewcommand{\arraystretch}{1.8}
        \fontsize{15}{21}\selectfont
        % Columna 1: Alineada a la derecha, en cursiva (itshape) y de 4.5cm de ancho
        % Columna 2: Alineada a la izquierda
        \begin{tabular}{>{\raggedleft\arraybackslash\itshape}p{4.5cm} @{\hspace{0.5cm}} l}
            Profesor(a): & \underline{\makebox[10cm][l]{ Eduardo Fragoso Navarro}} \\
            Asignatura: & \underline{\makebox[10cm][l]{ Fundamentos de Sistemas Embebidos}} \\
            Grupo: & \underline{\makebox[10cm][l]{3}} \\
            No de Práctica(s): & \underline{\makebox[10cm][l]{6}} \\
            Integrante(s): & \underline{\makebox[10cm][l]{Tepal Briseño Hansel Yael - 320078282}} \\
            & \underline{\makebox[10cm][l]{Ugartechea González Luis Antonio - 320223822}} \\
            No. de brigada: & \underline{\makebox[10cm][l]{12}} \\
            Semestre: & \underline{\makebox[10cm][l]{ 9\(^\circ\)}} \\
            Fecha de entrega: & \underline{\makebox[10cm][l]{13 de octubre del 2026}} \\
            Observaciones: & \underline{\makebox[10cm][l]{}} \\
            & \underline{\makebox[10cm][l]{}} \\
        \end{tabular}
    \end{center}

\end{titlepage}
\restoregeometry

\newpage
\fontsize{25}{21}\selectfont CALIFICACIÓN: \underline{\makebox[10cm][l]{}}
\newpage
\normalsize

\section*{\centering \fontsize{14}{17}\sffamily\selectfont Reconocimiento de Voz y de la Mano con IA}

\section*{\centering \fontsize{14}{17}\sffamily\selectfont Objetivos}

El objetivo de esta práctica fue integrar un sistema de captura distribuido en
una tarjeta XIAO ESP32-S3 con dos aplicaciones de inteligencia artificial
ejecutadas localmente en una computadora. En el primer ejercicio se recibió el
audio del micrófono mediante UDP y se transcribió al español con
\texttt{faster-whisper}. En el segundo se recibieron imágenes de la cámara y se
detectaron los puntos de una mano mediante MediaPipe Hand Landmarker.

También se buscó comprender la separación de responsabilidades entre el
dispositivo embebido y la computadora. El ESP32 realiza la adquisición, la
detección básica de actividad y el envío de datos, mientras que la PC ejecuta
los modelos de reconocimiento, que requieren más memoria y capacidad de
cómputo.

\section*{\centering \fontsize{14}{17}\sffamily\selectfont Resultados y Análisis}

\subsection*{\fontsize{12}{17}\sffamily\selectfont Configuración del sistema}

La comunicación se realizó en modo estación (STA). El ESP32 se conectó a la red
WiFi de 2.4 GHz y utilizó la dirección fija \texttt{192.168.8.212},
correspondiente al equipo 12. La computadora tuvo que estar conectada a la
misma red para poder intercambiar datagramas UDP con el puerto \texttt{3333}.
Se utilizó la misma estructura de firmware para los dos ejercicios y se
seleccionó la aplicación activa desde \texttt{main.c}.

El proyecto se compiló para el target \texttt{esp32s3}. El binario de audio
ocupó aproximadamente 751 kB y el binario de video 824 kB; ambos caben en la
partición de aplicación de 1.5 MB. El proceso de carga verificó los hashes de
los tres elementos escritos en la memoria flash y terminó con un reinicio
automático de la tarjeta.

La práctica activa se seleccionó desde el archivo principal. Solo una de las
dos funciones debe quedar definida para evitar ejecutar simultáneamente los dos
periféricos:

\begin{lstlisting}[style=cstyle, caption={Selección del ejercicio activo}, label={lst:main-practica}]
// Audio o video: dejar solamente una linea activa
// #define PRACTICA sta_audio_ahorro_app_main
#define PRACTICA sta_video_ahorro_app_main

void app_main(void)
{
	PRACTICA();
}
\end{lstlisting}

\subsection*{\fontsize{12}{17}\selectfont Ejercicio 2: voz a texto}

El programa del ESP32 muestrea el micrófono PDM a 16 kHz con muestras de 16
bits. Para reducir el consumo, el vigía permanece la mayor parte del tiempo en
reposo y solo enciende el micrófono periódicamente. Cuando detecta ruido por
tres mediciones consecutivas, comienza a enviar los bloques de audio que
superan el umbral configurado. La computadora se registra mediante un datagrama
\texttt{R} y recibe los bloques por UDP.

En \texttt{voz\_a\_texto.py} se acumulan 3 segundos de audio antes de llamar al
modelo. Esto representa:

\[
    3\,\mathrm{s}\times 16000\,\frac{\mathrm{muestras}}{\mathrm{s}}
    \times 2\,\frac{\mathrm{bytes}}{\mathrm{muestra}}=96000\,\mathrm{bytes}.
\]

Se utilizó el modelo \texttt{base} de \texttt{faster-whisper}, ejecutado en CPU
con \texttt{compute\_type="int8"}. Antes de transcribir, las muestras enteras
se convierten a valores de punto flotante en el intervalo aproximado $[-1,1]$.
La opción \texttt{vad\_filter=True} permite ignorar los segmentos de silencio y
el idioma se fija en español.

El bloque de transcripción empleado fue equivalente al siguiente:

\begin{lstlisting}[style=pythonstyle, caption={Conversión y transcripción del audio}, label={lst:whisper}]
audio = np.frombuffer(bloque, dtype=np.int16).astype(np.float32)
audio /= 32768.0
segmentos, _ = modelo.transcribe(
	audio, language="es", vad_filter=True)
for segmento in segmentos:
	print(segmento.text.strip(), flush=True)
\end{lstlisting}

El modelo \texttt{tiny} tendría menor latencia y consumo de memoria, pero
generalmente produciría más errores. \texttt{small} podría mejorar el
reconocimiento, especialmente con ruido o palabras menos frecuentes, a cambio
de mayor tiempo de procesamiento. El modelo \texttt{base} representa un
compromiso razonable para una computadora de propósito general.

El uso de bloques de 3 segundos introduce una latencia mínima de ese orden. Si
una palabra queda dividida entre dos bloques, puede perderse parte del contexto
y la transcripción puede ser incorrecta o quedar incompleta. Una mejora sería
usar bloques superpuestos o conservar una pequeña ventana del bloque anterior.

En el firmware, la tarea de audio solo envía un bloque cuando el nivel supera
el umbral. La dirección de la computadora se aprende cuando el script envía el
datagrama de registro:

\begin{lstlisting}[style=cstyle, caption={Registro de la PC y envío de audio por UDP}, label={lst:audio-udp}]
recibidos = recvfrom(s_socket, mensaje, sizeof(mensaje) - 1,
					 0, (struct sockaddr *)&origen,
					 &largo_origen);
if (recibidos > 0) {
	s_pc_direccion = origen;
	s_pc_conocida = true;
}

if (s_pc_conocida && microfono_nivel() > NIVEL_UMBRAL) {
	sendto(s_socket, s_bloque, leidos, 0,
		   (struct sockaddr *)&s_pc_direccion,
		   sizeof(s_pc_direccion));
}
\end{lstlisting}

\subsection*{\fontsize{12}{17}\selectfont Ejercicio 3: detección de la mano}

El firmware configura la cámara en escala de grises con una resolución de
$96\times96$ píxeles. El programa de la PC convierte cada imagen recibida de
BGR a RGB antes de entregarla a MediaPipe. El modelo Hand Landmarker se carga
una sola vez desde \texttt{hand\_landmarker.task} y se configura para buscar
una mano.

Para cada mano detectada se dibujan sus 21 puntos y las conexiones definidas
por MediaPipe. Además, se toman los puntos 4 y 8, correspondientes a la punta
del pulgar y del índice, respectivamente. La distancia mostrada se calcula en
píxeles mediante:

\[
    d=\sqrt{(x_8-x_4)^2+(y_8-y_4)^2}.
\]

La detección es más confiable cuando existe iluminación uniforme, la mano se
encuentra dentro del encuadre y el fondo contrasta con ella. Puede fallar con
poca luz, movimiento excesivo, oclusiones, fondos muy similares a la piel,
distancias demasiado grandes o cuando aparecen varias manos, ya que el programa
se configuró con \texttt{num\_hands=1}. La cámara permanece apagada durante los
intervalos de reposo; por ello la ventana de Python puede mostrar una imagen
negra mientras el ESP32 no ha detectado movimiento o la PC no está en la misma
red.

El procesamiento principal de cada imagen se realiza con la API Tasks de
MediaPipe. Primero se convierte la imagen recibida a RGB, luego se ejecuta el
detector y finalmente se dibujan las conexiones y la línea solicitada:

\begin{lstlisting}[style=pythonstyle, caption={Detección y dibujo de los puntos de la mano}, label={lst:mediapipe}]
rgb = cv2.cvtColor(gris, cv2.COLOR_GRAY2RGB)
imagen = mp.Image(image_format=mp.ImageFormat.SRGB, data=rgb)
resultado = detector.detect(imagen)

foto = cv2.cvtColor(gris, cv2.COLOR_GRAY2BGR)
for mano in resultado.hand_landmarks:
	dibujar_mano(foto, mano)

pulgar = puntos[4]
indice = puntos[8]
distancia = math.hypot(indice[0] - pulgar[0],
					   indice[1] - pulgar[1])
cv2.line(foto, pulgar, indice, (0, 255, 255), 2)
\end{lstlisting}

En el ESP32 la cámara se inicializa en escala de grises y con dos buffers para
mantener una imagen reciente disponible durante la transmisión:

\begin{lstlisting}[style=cstyle, caption={Configuración de la cámara del ESP32}, label={lst:camera-config}]
camera_config_t cam_cfg = {
	.pixel_format = PIXFORMAT_GRAYSCALE,
	.frame_size = FRAMESIZE_96X96,
	.fb_count = 2,
	.fb_location = CAMERA_FB_IN_DRAM,
	.grab_mode = CAMERA_GRAB_LATEST,
};
ESP_ERROR_CHECK(esp_camera_init(&cam_cfg));
\end{lstlisting}

\subsection*{\fontsize{12}{17}\selectfont Evidencias y problemas observados}

El proceso de carga produjo una verificación correcta de los datos escritos en
la memoria flash. Durante la prueba de video se observó inicialmente la ventana
negra de espera. La causa no fue MediaPipe, sino que la computadora no tenía
una ruta hacia \texttt{192.168.8.212}; estaba conectada a una red con dirección
\texttt{10.x.x.x}. Al conectar la PC a la misma red WiFi del ESP32, el script
pudo registrarse mediante UDP.

Para completar la evidencia solicitada por la práctica se deben agregar al
directorio \texttt{reporte/img/} una captura de la terminal con una
transcripción y una captura de la ventana de video con los puntos de la mano, y
después insertarlas en esta sección. No se incluyeron imágenes inventadas
porque dichas capturas no estaban disponibles en el directorio del reporte.

\section*{\centering \fontsize{14}{17}\sffamily\selectfont Conclusiones}

El reconocimiento se ejecuta en la PC porque los modelos de Whisper y MediaPipe
requieren más memoria, almacenamiento y capacidad de cómputo que la disponible
en el ESP32 para esta implementación. Para ejecutar el reconocimiento
directamente en el microcontrolador sería necesario usar modelos mucho más
pequeños, cuantizados y optimizados, además de disponer de memoria suficiente y
aceptar una menor precisión.

El procesamiento local tiene como ventajas que no requiere enviar audio o video
a un servicio externo, puede funcionar sin Internet y ofrece mayor control
sobre la privacidad. Sus desventajas son la necesidad de instalar los modelos,
ocupar espacio en la computadora y disponer de suficiente capacidad de
procesamiento. Un servicio en la nube simplifica la ejecución y puede ofrecer
modelos más grandes, pero depende de Internet, agrega latencia y expone los
datos capturados a un tercero.

La arquitectura puede utilizarse en asistentes de voz locales, sistemas de
accesibilidad, control de dispositivos mediante palabras, interfaces sin
contacto, interpretación de señas, monitoreo de maquinaria y automatización
doméstica. Para hacerla más inteligente se podrían agregar clasificación de
palabras clave, identificación de hablantes, reducción de ruido, seguimiento
temporal de la mano y un clasificador de gestos entrenado con los 21 puntos.
Así sería posible reconocer comandos completos o traducir un conjunto de señas
a acciones concretas.

% \section{Bibliografía}

\begin{biblio}
	% for: https://sistemasmultibasededatos.blogspot.com/
	\item Esquivel Martinez H. (2021, Noviembre 28). \textit{SISTEMAS MULTIBASE DE DATOS} [Online]. Disponible en: \newline \textcolor{blue}{\underline{\url{https://sistemasmultibasededatos.blogspot.com/}}}
\end{biblio}

\end{document}
